% ===========================================================================
% preamble.tex — shared packages, styles and macros for the bitterfs docs
% ===========================================================================
%
% Engine: lualatex.  Chosen for native UTF-8 and real system fonts, because
% these documents are full of em-dashes, arrows and box-drawing characters that
% pdflatex would need escaping for.
%
% Build with:  make -C docs
% ===========================================================================

% --- Page and fonts --------------------------------------------------------
\usepackage{geometry}
\geometry{a4paper, margin=2.5cm}

% fontspec needs lualatex or xelatex.  Latin Modern is the TeX default and is
% always present; DejaVu Sans Mono is chosen for code specifically because of
% its wide glyph coverage — box-drawing characters and arrows appear inside
% verbatim blocks throughout these documents.
\usepackage{fontspec}
\setmainfont{Latin Modern Roman}
\setsansfont{DejaVu Sans}[Scale=MatchLowercase]
\setmonofont{DejaVu Sans Mono}[Scale=MatchLowercase]

\usepackage{microtype}

% --- Tables ----------------------------------------------------------------
\usepackage{booktabs}   % \toprule/\midrule/\bottomrule — no vertical rules
\usepackage{longtable}  % tables that break across pages (pandoc emits these)
\usepackage{array}
\usepackage{multirow}

% --- Lists, colour, captions -----------------------------------------------
\usepackage{enumitem}
\usepackage{xcolor}
\usepackage[font=small,labelfont=bf]{caption}

% --- Diagrams --------------------------------------------------------------
% bytefield draws on-disk structure layouts with byte/bit offset rulers.  This
% is the package that makes LaTeX worth the trouble for a format specification:
% \wordbox and \bitbox produce a real diagram of a struct rather than ASCII art.
\usepackage{bytefield}

% tikz for everything bytefield does not cover — chiefly the leaf layout, where
% items grow forward from the header while payloads stack backward from the end
% of the block.
\usepackage{tikz}
\usetikzlibrary{positioning, arrows.meta, fit, backgrounds, calc, shapes.multipart}

% --- Code listings ---------------------------------------------------------
\usepackage{listings}

\definecolor{codekw}{HTML}{0B5FA5}
\definecolor{codecm}{HTML}{6A737D}
\definecolor{codestr}{HTML}{A31515}
\definecolor{codeframe}{HTML}{C8C8C8}
\definecolor{codebg}{HTML}{FAFAFA}

% listings processes input character by character, so multi-byte UTF-8 breaks
% it even under lualatex.  The `literate` table below maps each character that
% actually appears in these documents to something typesettable.  Add to it
% rather than removing the character from the source — the source is also read
% as plain text in a terminal, where the box drawing is the point.
\lstset{
  basicstyle=\ttfamily\small,
  keywordstyle=\color{codekw}\bfseries,
  commentstyle=\color{codecm}\itshape,
  stringstyle=\color{codestr},
  numberstyle=\tiny\color{codecm},
  backgroundcolor=\color{codebg},
  frame=single,
  rulecolor=\color{codeframe},
  framesep=4pt,
  breaklines=true,
  breakatwhitespace=true,
  showstringspaces=false,
  columns=fullflexible,
  keepspaces=true,
  upquote=true,
  extendedchars=true,
  literate=
    {─}{{-}}1 {│}{{|}}1 {┌}{{+}}1 {┐}{{+}}1
    {└}{{+}}1 {┘}{{+}}1 {├}{{+}}1 {┤}{{+}}1 {┬}{{+}}1 {┴}{{+}}1 {┼}{{+}}1
    {→}{{$\rightarrow$}}1 {←}{{$\leftarrow$}}1
    {↑}{{$\uparrow$}}1    {↓}{{$\downarrow$}}1
    {≤}{{$\leq$}}1        {≥}{{$\geq$}}1  {≈}{{$\approx$}}1
    {×}{{$\times$}}1      {—}{{---}}1     {–}{{--}}1
    {“}{{``}}1            {”}{{''}}1      {‘}{{`}}1  {’}{{'}}1
    {…}{{\ldots}}1        {·}{{$\cdot$}}1
}

% Styles are DEFINED here but no global default language is set.  An earlier
% version made C the default, which meant every ASCII diagram containing the
% word "struct" got it bold-blued.  Ask for the style you want:
%     \begin{lstlisting}[style=c]
\lstdefinestyle{c}{language=C}

% Shell transcripts and hexdumps: no keyword highlighting, since C keywords
% appearing in a shell command would be coloured misleadingly.
\lstdefinestyle{shell}{language=,keywordstyle=,basicstyle=\ttfamily\small}

% --- Code blocks in CONVERTED documents ------------------------------------
% pandoc emits Shaded/Highlighting rather than lstlisting.  That is the better
% path for these documents and the reason listings is not used globally:
% Highlighting is built on fancyvrb, which under lualatex passes UTF-8 straight
% through, so the box-drawing characters in the architecture diagrams survive.
% listings processes input one character at a time and mangles them.
\usepackage{framed}
\definecolor{shadecolor}{HTML}{FAFAFA}
% ===========================================================================
% pandoc-highlight.tex — syntax-highlighting definitions required by pandoc
% ===========================================================================
%
% GENERATED, do not hand-edit.  Regenerate with:  make -C docs highlight
%
% pandoc's LaTeX writer wraps code blocks in \begin{Shaded}\begin{Highlighting}
% and colours them with one macro per token class (\KeywordTok, \CommentTok,
% ...).  None of that is defined by any package, so a converted document does
% not build without these definitions.
%
% This is why listings is not used for converted documents: Highlighting is
% built on fancyvrb, which under lualatex passes UTF-8 through natively, so the
% box-drawing characters in the architecture diagrams survive.  listings
% processes input one character at a time and mangles them.
% ===========================================================================

\usepackage{color}
\usepackage{fancyvrb}
\newcommand{\VerbBar}{|}
\newcommand{\VERB}{\Verb[commandchars=\\\{\}]}
\DefineVerbatimEnvironment{Highlighting}{Verbatim}{commandchars=\\\{\}}
% Add ',fontsize=\small' for more characters per line
\newenvironment{Shaded}{}{}
\newcommand{\AlertTok}[1]{\textcolor[rgb]{1.00,0.00,0.00}{\textbf{#1}}}
\newcommand{\AnnotationTok}[1]{\textcolor[rgb]{0.38,0.63,0.69}{\textbf{\textit{#1}}}}
\newcommand{\AttributeTok}[1]{\textcolor[rgb]{0.49,0.56,0.16}{#1}}
\newcommand{\BaseNTok}[1]{\textcolor[rgb]{0.25,0.63,0.44}{#1}}
\newcommand{\BuiltInTok}[1]{\textcolor[rgb]{0.00,0.50,0.00}{#1}}
\newcommand{\CharTok}[1]{\textcolor[rgb]{0.25,0.44,0.63}{#1}}
\newcommand{\CommentTok}[1]{\textcolor[rgb]{0.38,0.63,0.69}{\textit{#1}}}
\newcommand{\CommentVarTok}[1]{\textcolor[rgb]{0.38,0.63,0.69}{\textbf{\textit{#1}}}}
\newcommand{\ConstantTok}[1]{\textcolor[rgb]{0.53,0.00,0.00}{#1}}
\newcommand{\ControlFlowTok}[1]{\textcolor[rgb]{0.00,0.44,0.13}{\textbf{#1}}}
\newcommand{\DataTypeTok}[1]{\textcolor[rgb]{0.56,0.13,0.00}{#1}}
\newcommand{\DecValTok}[1]{\textcolor[rgb]{0.25,0.63,0.44}{#1}}
\newcommand{\DocumentationTok}[1]{\textcolor[rgb]{0.73,0.13,0.13}{\textit{#1}}}
\newcommand{\ErrorTok}[1]{\textcolor[rgb]{1.00,0.00,0.00}{\textbf{#1}}}
\newcommand{\ExtensionTok}[1]{#1}
\newcommand{\FloatTok}[1]{\textcolor[rgb]{0.25,0.63,0.44}{#1}}
\newcommand{\FunctionTok}[1]{\textcolor[rgb]{0.02,0.16,0.49}{#1}}
\newcommand{\ImportTok}[1]{\textcolor[rgb]{0.00,0.50,0.00}{\textbf{#1}}}
\newcommand{\InformationTok}[1]{\textcolor[rgb]{0.38,0.63,0.69}{\textbf{\textit{#1}}}}
\newcommand{\KeywordTok}[1]{\textcolor[rgb]{0.00,0.44,0.13}{\textbf{#1}}}
\newcommand{\NormalTok}[1]{#1}
\newcommand{\OperatorTok}[1]{\textcolor[rgb]{0.40,0.40,0.40}{#1}}
\newcommand{\OtherTok}[1]{\textcolor[rgb]{0.00,0.44,0.13}{#1}}
\newcommand{\PreprocessorTok}[1]{\textcolor[rgb]{0.74,0.48,0.00}{#1}}
\newcommand{\RegionMarkerTok}[1]{#1}
\newcommand{\SpecialCharTok}[1]{\textcolor[rgb]{0.25,0.44,0.63}{#1}}
\newcommand{\SpecialStringTok}[1]{\textcolor[rgb]{0.73,0.40,0.53}{#1}}
\newcommand{\StringTok}[1]{\textcolor[rgb]{0.25,0.44,0.63}{#1}}
\newcommand{\VariableTok}[1]{\textcolor[rgb]{0.10,0.09,0.49}{#1}}
\newcommand{\VerbatimStringTok}[1]{\textcolor[rgb]{0.25,0.44,0.63}{#1}}
\newcommand{\WarningTok}[1]{\textcolor[rgb]{0.38,0.63,0.69}{\textbf{\textit{#1}}}}

% fvextra extends fancyvrb with line breaking; plain fancyvrb has no
% breaklines option at all, and asking for one is a keyval error.
\usepackage{fvextra}

% fancyvrb does not wrap by default, so long lines run off the page.
% breakanywhere is required because the offenders are paths and hexdumps with
% no spaces to break at.
\RecustomVerbatimEnvironment{Highlighting}{Verbatim}{%
  commandchars=\\\{\},
  breaklines=true, breakanywhere=true,
  fontsize=\small}

\renewenvironment{Shaded}{\begin{snugshade}}{\end{snugshade}}

% --- pandoc compatibility --------------------------------------------------
% pandoc's LaTeX writer emits \tightlist for lists without blank lines between
% items, and \passthrough around inline code in some contexts.  Neither is
% defined by any package, so converted documents fail to build without them.
\providecommand{\tightlist}{%
  \setlength{\itemsep}{0pt}\setlength{\parskip}{0pt}}
\providecommand{\passthrough}[1]{#1}

% --- Cross-references ------------------------------------------------------
% hyperref must come late; cleveref must come after hyperref.
\usepackage[hidelinks, bookmarksnumbered]{hyperref}
\usepackage{cleveref}

\usepackage[backend=biber, style=numeric, sorting=none]{biblatex}
\addbibresource{refs.bib}


% --- Project macros --------------------------------------------------------
% Consistent typesetting for the things these documents name constantly.  Using
% a macro rather than \texttt directly means the styling is changeable in one
% place, and it documents intent: \key{} is a key, not just monospace text.
\newcommand{\code}[1]{\texttt{#1}}
\newcommand{\fnname}[1]{\texttt{#1}}                       % a function
\newcommand{\field}[1]{\texttt{#1}}                        % a struct field
\newcommand{\itemtype}[1]{\textsc{\MakeLowercase{#1}}}     % INODE_ITEM etc.

% A (objectid, type, offset) key, typeset consistently.
% \mbox rather than \text: \text needs amsmath, which is not loaded, and this
% macro went unused for long enough that nobody noticed it could not compile.
\newcommand{\keytriple}[3]{\ensuremath{(\mbox{\code{#1}},\,\mbox{\code{#2}},\,\mbox{\code{#3}})}}

% Byte-count shorthand used throughout the format spec.
\newcommand{\bytes}[1]{#1\,B}
\newcommand{\kib}[1]{#1\,KiB}
\newcommand{\mib}[1]{#1\,MiB}

% An open decision, called out visibly in the rendered PDF.  A format spec with
% unresolved choices buried in prose is a format spec you will implement twice;
% these boxes are meant to be impossible to skim past, and to disappear as they
% are settled.
\newcommand{\decide}[1]{%
  \par\medskip\noindent
  \fcolorbox{black!30}{yellow!12}{%
    \parbox{\dimexpr\linewidth-2\fboxsep-2\fboxrule}{%
      \textbf{Decide:}\ #1}}%
  \par\medskip}
